\documentclass[11pt,twocolumn,a4paper]{article}

% ─── Packages ─────────────────────────────────────────────────────────
\usepackage[margin=2.0cm,top=2.4cm,bottom=2.4cm]{geometry}
\usepackage[utf8]{inputenc}
\usepackage[T1]{fontenc}
\usepackage{lmodern}
\usepackage{microtype}
\usepackage{amsmath,amssymb,amsthm,mathtools}
\usepackage{physics}
\usepackage{bm}
\usepackage{graphicx}
\usepackage{booktabs}
\usepackage{array}
\usepackage{multirow}
\usepackage{enumitem}
\usepackage{float}
\usepackage{caption}
\usepackage{natbib}
\usepackage{xcolor}
\usepackage[colorlinks=true,linkcolor=blue!60!black,
            citecolor=blue!60!black,urlcolor=blue!60!black]{hyperref}
\usepackage{fancyhdr}
\usepackage{titlesec}

% ─── Page Style ───────────────────────────────────────────────────────
\pagestyle{fancy}
\fancyhf{}
\fancyhead[L]{\small\textit{Ion Channel Aperture Arrays as Categorical Filters}}
\fancyhead[R]{\small\thepage}
\renewcommand{\headrulewidth}{0.4pt}

\titleformat{\section}{\large\bfseries}{\thesection.}{0.5em}{}
\titleformat{\subsection}{\normalsize\bfseries}{\thesubsection.}{0.5em}{}
\titleformat{\subsubsection}{\normalsize\itshape}{\thesubsubsection.}{0.5em}{}

% ─── Theorem Environments ─────────────────────────────────────────────
\newtheorem{theorem}{Theorem}[section]
\newtheorem{lemma}[theorem]{Lemma}
\newtheorem{corollary}[theorem]{Corollary}
\newtheorem{proposition}[theorem]{Proposition}
\theoremstyle{definition}
\newtheorem{definition}[theorem]{Definition}
\newtheorem{axiom}{Axiom}
\theoremstyle{remark}
\newtheorem{remark}[theorem]{Remark}

% ─── Custom Commands ──────────────────────────────────────────────────
\newcommand{\kB}{k_{\mathrm{B}}}
\newcommand{\R}{\mathbb{R}}
\newcommand{\N}{\mathbb{N}}
\newcommand{\Sk}{S_{\mathrm{k}}}
\newcommand{\St}{S_{\mathrm{t}}}
\newcommand{\Se}{S_{\mathrm{e}}}
\newcommand{\kap}{\kappa}
\newcommand{\Sent}{\bm{S}}

% ─── Caption Setup ────────────────────────────────────────────────────
\captionsetup{font=small,labelfont=bf}

% ══════════════════════════════════════════════════════════════════════
\begin{document}

\twocolumn[
\begin{@twocolumnfalse}
\begin{center}

{\LARGE\bfseries Ion Channel Aperture Arrays as Categorical Filters:\\[4pt]
The Five Names Theorem, Catalytic Power, and Saturation Dynamics}

\vspace{12pt}

{\large Kundai Farai Sachikonye}

\vspace{6pt}

{\small\itshape Chair of Experimental Bioinformatics, Technical University of Munich \\[0.3em]
\small\itshape AIMe Registry for Artificial Intelligence \\[0.2em]
\small\itshape Maximus-von-Imhof-Forum 3, 85354 Freising, Germany \\[0.2em]
\texttt{kundai.sachikonye@wzw.tum.de}}

\vspace{6pt}{\small May 2026}

\vspace{16pt}

\begin{minipage}{0.92\textwidth}
\textbf{Abstract.}
We derive ion channel aperture arrays as mathematical necessities of bounded
partition-based computation in aqueous lipid-bilayer substrates. The argument
proceeds from the Bounded Phase Space axiom through five theorems to a unique
aperture implementation. A partition boundary separating two regions of phase
space with finite Liouville measure presents any crossing particle with a Born
desolvation energy barrier of approximately $40\,\kB T$ in a hydrophobic slab
of dielectric constant $\varepsilon_m \approx 2$; spontaneous crossing is
therefore exponentially suppressed ($\sim e^{-40} \approx 10^{-17}$) and
computation requires specialised permeation structures. We prove the Five Names
Theorem: any selective permeation structure at the partition boundary of a
bounded aqueous dynamical system is simultaneously a geometric aperture with
transmission function $\kap \in \lbrack 0,1\rbrack$, a mediator of virtual sub-states in
the backward trajectory completion algorithm, an information catalyst for the
crossing transformation, a Maxwell demon for S-entropy coordinate measurement,
and a physically instantiated subtask node in the cascade-fabric routing
graph. These five characterisations are shown to be pairwise equivalent.
In an aqueous lipid bilayer at physiological temperature, the unique physical
realisation satisfying all five is an ion-channel protein: a water-filled
protein pore that collapses the Born barrier, provides a selectivity filter
(S-entropy matching gate), and is structurally invariant across individual ion
transit events. The catalytic power $\kap(\mathcal{X}; \mathbf{t})$ of a
channel with target S-entropy coordinates $\mathbf{t}$ for an input system
with coordinates $\Sent(\mathcal{X})$ equals the spectral overlap integral
between their Koopman spectra. For an array of $n$ channels in parallel, the
array catalytic power satisfies $\kap_{\mathrm{arr}} = 1 - \prod_{i=1}^n (1-\kap_i)$
(the Saturation Theorem), which converges to unity if and only if
$\sum_{i=1}^n \kap_i = \infty$. The action potential is identified as a
virtual sub-state: a finite-lifetime Koopman excitation that enables the
membrane to transition between conductance eigenstates without permanently
occupying either. Validation against patch-clamp conductance data confirms
predicted catalytic powers to within experimental uncertainty.
\end{minipage}

\vspace{8pt}
\noindent\textbf{Keywords:}
ion channel, Born energy barrier, Maxwell demon, catalytic power, Five Names Theorem,
virtual sub-state, action potential, cascade saturation, bounded phase space,
aperture array

\vspace{18pt}
\end{center}
\end{@twocolumnfalse}
]

%──────────────────────────────────────────────────────────────────────
\section{Introduction}
\label{sec:introduction}
%──────────────────────────────────────────────────────────────────────

The lipid bilayer membrane is, as established in the first companion paper
\citep{companion1}, the necessary and sufficient partition boundary for
computation in bounded aqueous systems. A partition boundary that is completely
impermeable separates two phase space regions permanently: no information can
cross, and no computation is possible. A boundary that is fully permeable
cannot sustain a partition: the two regions merge into one. Useful computation
requires a boundary that is \emph{selectively} permeable --- one that passes
some trajectories and blocks others according to a computable criterion.

The question of how selective permeability arises in a lipid bilayer has been
answered in biology by five decades of electrophysiology: ion channels are
specialised protein structures that create water-filled pores spanning the
membrane, equipped with selectivity filters and voltage- or ligand-controlled
gates \citep{hille2001, neher1976, doyle1998}. The standard account is
evolutionary: channels evolved over billions of years to satisfy physiological
requirements. This account is correct but incomplete: it does not explain
\emph{why} the resulting structures have the specific mathematical properties
they do, nor why those properties are precisely what is needed for the partition
computation described in this paper series.

We give a first-principles derivation. Starting from the Born energy barrier
for ion crossing of a hydrophobic slab \citep{parsegian1969}, we show that
the bilayer is exponentially impermeable to free ions
(\S\ref{sec:opacity}). We then derive, from the operational requirements of
partition-based computation, the set of properties that any permeation structure
must satisfy (\S\ref{sec:fivenames}). The Five Names Theorem shows that these
properties are not five independent requirements but five descriptions of a
single mathematical object. In an aqueous bilayer, this object is uniquely
realised by an ion channel (\S\ref{sec:channelnec}). The catalytic power
$\kap$ --- the probability that an input system with given S-entropy coordinates
passes through a given aperture --- is derived from the Koopman spectral overlap
(\S\ref{sec:catpower}). The action potential is identified as a virtual sub-state
mediating transitions between conductance eigenstates (\S\ref{sec:virtsubstate}).
The Saturation Theorem for parallel arrays is proved in \S\ref{sec:saturation}.

%──────────────────────────────────────────────────────────────────────
\section{Bounded Phase Space Framework}
\label{sec:bps}
%──────────────────────────────────────────────────────────────────────

We state the foundational axiom and collect the results from companion papers
\citep{companion1, companion2} needed below.

\begin{axiom}[Bounded Phase Space]
\label{ax:bps}
Every persistent physical system $\mathcal{X}$ occupies a region
$\Omega \subset \mathbb{R}^{2d}$ with finite Liouville measure
$\mu(\Omega) = \int_\Omega d^d q\, d^d p / h^d < \infty$,
admitting a hierarchy of refining measurable partitions
$\mathscr{P}_0 \supset \mathscr{P}_1 \supset \cdots$.
\end{axiom}

\noindent From \citep{companion1}: The partition boundary $\partial\Omega$
separating interior $\Omega_{\mathrm{in}}$ from exterior $\Omega_{\mathrm{out}}$
is a lipid bilayer membrane of thickness $d = 4.0$~nm and bending modulus
$\kap_b = 20\,\kB T$.

\noindent From \citep{companion2}: Every bounded oscillatory system has a
discrete Koopman spectrum $\Sigma(\mathcal{X}) = \{(\omega_k, A_k, \phi_k)\}$,
and its S-entropy coordinates $\Sent(\mathcal{X}) = (\Sk, \St, \Se) \in [0,1]^3$
are the unique minimal sufficient statistics for categorical identification.

%──────────────────────────────────────────────────────────────────────
\section{Partition Boundary Opacity}
\label{sec:opacity}
%──────────────────────────────────────────────────────────────────────

\subsection{The Born Energy Barrier}

A monovalent ion of radius $r$ and charge $Ze$ crossing from water
($\varepsilon_w \approx 80$) into a hydrophobic slab ($\varepsilon_m \approx 2$)
of thickness $d$ experiences a desolvation energy barrier. The Born
model \citep{born1920} gives, for an infinite slab, the midpoint energy:
\begin{equation}
\Delta E_{\mathrm{Born}} = \frac{(Ze)^2}{8\pi\varepsilon_0 r}
\left(\frac{1}{\varepsilon_m} - \frac{1}{\varepsilon_w}\right).
\label{eq:born}
\end{equation}
For a slab of finite thickness $d$, the image-charge series
gives a correction; the full result for $d = 4$~nm,
$r = 0.095$~nm (Na$^+$) was computed by \citet{parsegian1969}:
\begin{equation}
\Delta E_{\mathrm{Born}}^{(\mathrm{slab})} \approx 40\,\kB T \quad
(T = 310\,\mathrm{K}).
\label{eq:born_slab}
\end{equation}

\begin{theorem}[Partition Boundary Opacity]
\label{thm:opacity}
The spontaneous crossing rate of monovalent ions across a lipid bilayer
membrane in the absence of specialised permeation structures is:
\begin{equation}
k_{\mathrm{cross}} = \nu_0 \exp(-\Delta E_{\mathrm{Born}} / \kB T)
\approx 10^{9}\,\mathrm{s}^{-1} \times e^{-40}
\approx 4 \times 10^{-9}\,\mathrm{s}^{-1}\,\mathrm{ion}^{-1},
\label{eq:crossrate}
\end{equation}
where $\nu_0 \approx \kB T / h \approx 6 \times 10^{12}\,\mathrm{s}^{-1}$ is
the attempt frequency (thermal rate at the transition state).
The corresponding permeability coefficient is
$P_{\mathrm{ion}} \approx 10^{-14}$~cm~s$^{-1}$, seventeen orders of magnitude
below the water permeability coefficient $P_{\mathrm{water}} \approx 10^{3}$~cm~s$^{-1}$.
\end{theorem}

\begin{proof}
Equation~\eqref{eq:crossrate} follows directly from
equation~\eqref{eq:born_slab} and Kramers' transition-state rate theory
\citep{kramers1940} with $\nu_0 = \kB T / h$ (the thermal frequency).
The permeability coefficient $P = k_{\mathrm{cross}} \times d = 4 \times
10^{-9}\,\mathrm{s}^{-1} \times 4 \times 10^{-9}\,\mathrm{m} \approx
10^{-17}\,\mathrm{m\,s}^{-1} = 10^{-15}\,\mathrm{cm\,s}^{-1}$, consistent
with the experimental lower bound of $P_{\mathrm{Na}^+} \leq 10^{-14}$~cm~s$^{-1}$
measured for pure bilayers \citep{hille2001, walter1983}.
\end{proof}

\begin{corollary}[Computational Stasis]
\label{cor:stasis}
A partition boundary with no specialised permeation structures supports
zero information flux between $\Omega_{\mathrm{in}}$ and $\Omega_{\mathrm{out}}$
on any timescale relevant to partition-based computation ($\tau \lesssim 1$~s).
\end{corollary}

\begin{proof}
At the spontaneous crossing rate~\eqref{eq:crossrate}, the expected number of
spontaneous crossings per square micrometre of membrane per second is
$k_{\mathrm{cross}} \times \rho_{\mathrm{ion}} \times d \approx 10^{-17}$~m~s$^{-1}
\times 0.15 \times 10^3$~mol~L$^{-1} \times N_A \approx 10^{-6}$~ion~$\mu$m$^{-2}$~s$^{-1}$.
Since each crossing carries at most $\kB T \ln 2$ bits of information
\citep{landauer1961}, the information flux is $\lesssim 10^{-6} \kB T \ln 2
\,\mu\mathrm{m}^{-2}\,\mathrm{s}^{-1}$, negligible compared to the
$\sim 10^{20}$~ops~s$^{-1}$~cm$^{-2}$ capacity of the fluid-phase membrane
\citep{companion1}.
\end{proof}

Corollary~\ref{cor:stasis} establishes the necessity of specialised permeation
structures. We now derive the properties such structures must have.

\begin{figure*}[!htbp]
    \centering
    \includegraphics[width=\textwidth]{figures/aa_01.png}
    \caption{\textbf{Born Energy Barrier for Na$^+$ Crossing the Bilayer Core Is $\gg k_BT$.}
    (\textbf{A})~Bulk Born transfer energy $\Delta W_\mathrm{Born}(r) = (ze)^2/(8\pi\varepsilon_0 r)\,(1/\varepsilon_l - 1/\varepsilon_w)$ in units of $k_BT$ versus ionic radius $r \in [0.04, 0.50]$~\si{\nano\metre}, for $\varepsilon_l = 2.2$ and $\varepsilon_w = 80$. The vertical dashed line marks the Na$^+$ crystal radius $r = 0.095$~\si{\nano\metre}; the dotted horizontal at $40\,k_BT$ is the Parsegian refined estimate (Parsegian 1969; Hille 2001). The bulk Born barrier at $r = 0.095$~\si{\nano\metre} exceeds $100\,k_BT$, confirming partition opacity.
    (\textbf{B})~Born barrier versus lipid dielectric constant $\varepsilon_l \in [1.5, 10]$ at fixed $r = 0.095$~\si{\nano\metre}. The barrier decreases with $\varepsilon_l$ but remains above $40\,k_BT$ for all physically realistic lipid matrices, establishing the robustness of the opacity result.
    (\textbf{C})~Convergence of the Parsegian image-charge series $\sum_{n=1}^{N} \Delta^n/n$ versus the number of retained terms $N$. The analytic limit $-\ln(1-\Delta)$ (dashed) is attained by $N \approx 30$; the fast geometric convergence validates the use of the truncated series in claim aa\_01.
    (\textbf{D})~Three-dimensional surface of $\Delta W_\mathrm{Born}$ clipped to $400\,k_BT$ over the joint plane $(r, \varepsilon_l) \in [0.05, 0.4]$~\si{\nano\metre} $\times$ $[1.5, 10]$ (hot\_r palette). The divergence as $r \to 0$ and $\varepsilon_l \to 1$ makes the worst-case scenario visible; the physiological regime lies on the descending flank far above the $10\,k_BT$ opacity criterion.}\label{fig:aa_01}
\end{figure*}

%──────────────────────────────────────────────────────────────────────
\section{The Five Names Theorem}
\label{sec:fivenames}
%──────────────────────────────────────────────────────────────────────

\subsection{Five Definitions of an Aperture}

We give five definitions, each capturing a different aspect of a selective
permeation structure at the partition boundary.

\begin{definition}[Geometric Aperture]
\label{def:geom}
A \emph{geometric aperture} at the partition boundary $\partial\Omega$ is a
pair $(A, \kap)$ where $A \subseteq \partial\Omega$ is a measurable region
(the \emph{pore domain}) and $\kap: [0,1]^3 \to [0,1]$ is a
\emph{transmission function} mapping an input system's S-entropy coordinates
$\Sent(\mathcal{X})$ to the probability that it crosses $\partial\Omega$
through $A$.
\end{definition}

\begin{definition}[Virtual Sub-State Mediator]
\label{def:virtual}
A \emph{virtual sub-state mediator} is a structure $V$ at $\partial\Omega$
that enables passage of a trajectory from $\Omega_{\mathrm{in}}$ to
$\Omega_{\mathrm{out}}$ (or vice versa) via a transient state $|v\rangle$
satisfying:
\begin{enumerate}[nosep, label=(\roman*)]
\item $|v\rangle$ has non-zero Koopman amplitude: $A_v > 0$;
\item $|v\rangle$ is not an eigenstate of the partition operator $\hat{P}$
  of either $\Omega_{\mathrm{in}}$ or $\Omega_{\mathrm{out}}$:
  $\hat{P}|v\rangle \neq \lambda|v\rangle$;
\item $|v\rangle$ has finite lifetime: $\tau_v = A_v^{-1}(\mathrm{d}A_v/\mathrm{d}t)^{-1}
  < \infty$.
\end{enumerate}
\end{definition}

\begin{definition}[Information Catalyst]
\label{def:catalyst}
An \emph{information catalyst} for the crossing transformation
$T: \Omega_{\mathrm{in}} \to \Omega_{\mathrm{out}}$ is an agent $C$ satisfying:
\begin{enumerate}[nosep, label=(\roman*)]
\item \textit{Necessity}: $T$ has zero rate without $C$
  (Theorem~\ref{thm:opacity});
\item \textit{Catalytic invariance}: the structure of $C$ is unchanged by each
  transit event: $|C\rangle_{\mathrm{before}} = |C\rangle_{\mathrm{after}}$;
\item \textit{Information conservation}: $H(C_{\mathrm{before}}) = H(C_{\mathrm{after}})$,
  where $H$ is the Shannon entropy of the channel's conformational state
  distribution.
\end{enumerate}
\end{definition}

\begin{definition}[Maxwell Demon]
\label{def:demon}
A \emph{Maxwell demon} at the partition boundary is an agent $D$ that:
\begin{enumerate}[nosep, label=(\roman*)]
\item Measures the S-entropy coordinates $\Sent(\mathcal{X})$ of each approaching
  particle;
\item Conditionally permits crossing with probability $\kap(\Sent(\mathcal{X}))$
  based on the measurement outcome;
\item Operates at energy cost $\geq \kB T \ln 2$ per binary decision
  \citep{szilard1929, landauer1961}.
\end{enumerate}
\end{definition}

\begin{definition}[Cascade Subtask Node]
\label{def:subtask}
A \emph{cascade subtask node} in the backward trajectory completion algorithm
is a node $u$ in the routing graph $G = (V, E)$ of the cascade fabric such
that:
\begin{enumerate}[nosep, label=(\roman*)]
\item $u$ lies on the boundary between two ternary routing sub-trees;
\item Passage through $u$ is conditional: a trajectory routed to $u$
  is forwarded to the sub-tree on the other side with probability $\kap_u$
  and reflected with probability $1 - \kap_u$;
\item $\kap_u$ is computed from the S-entropy distance between the trajectory's
  coordinates and the sub-tree's target category.
\end{enumerate}
\end{definition}

\subsection{Equivalence}

\begin{theorem}[Five Names]
\label{thm:fivenames}
For any structure $\mathcal{A}$ at the partition boundary $\partial\Omega$ of a
bounded aqueous dynamical system, the following are equivalent:
\begin{itemize}[nosep]
\item[(\textbf{G})] $\mathcal{A}$ is a geometric aperture (Definition~\ref{def:geom});
\item[(\textbf{V})] $\mathcal{A}$ is a virtual sub-state mediator
  (Definition~\ref{def:virtual});
\item[(\textbf{I})] $\mathcal{A}$ is an information catalyst
  (Definition~\ref{def:catalyst});
\item[(\textbf{M})] $\mathcal{A}$ is a Maxwell demon (Definition~\ref{def:demon});
\item[(\textbf{S})] $\mathcal{A}$ is a cascade subtask node
  (Definition~\ref{def:subtask}).
\end{itemize}
\end{theorem}

\begin{proof}
We prove the cycle $\textbf{G} \Rightarrow \textbf{V} \Rightarrow \textbf{I}
\Rightarrow \textbf{M} \Rightarrow \textbf{S} \Rightarrow \textbf{G}$.

\vspace{4pt}
\noindent$\textbf{G} \Rightarrow \textbf{V}$: A geometric aperture with
$0 < \kap < 1$ has a pore domain $A \subseteq \partial\Omega$. A particle
traversing $A$ occupies $A$ transiently during transit, with duration
$\tau_v \approx L_{\mathrm{pore}} / v_{\mathrm{th}}$ where $L_{\mathrm{pore}}$
is the pore length and $v_{\mathrm{th}} = \sqrt{2\kB T / m}$ is the thermal
velocity. During this transit:

(i) The particle's trajectory has non-zero Koopman amplitude at frequencies
corresponding to the transit: $A_v = |F(\omega_{\mathrm{transit}})|^2 > 0$
since $\omega_{\mathrm{transit}} = v_{\mathrm{th}}/L_{\mathrm{pore}} > 0$.

(ii) The transit state is not an eigenstate of $\hat{P}_{\mathrm{in}}$ or
$\hat{P}_{\mathrm{out}}$ since the particle is in neither $\Omega_{\mathrm{in}}$
nor $\Omega_{\mathrm{out}}$ but on the boundary $\partial\Omega$.

(iii) The lifetime $\tau_v = L_{\mathrm{pore}}/v_{\mathrm{th}} < \infty$
since both $L_{\mathrm{pore}}$ and $v_{\mathrm{th}}$ are positive and finite.

Hence $\mathcal{A}$ mediates a virtual sub-state.

\vspace{4pt}
\noindent$\textbf{V} \Rightarrow \textbf{I}$: A virtual sub-state mediator $V$
enables a crossing event (from $\Omega_{\mathrm{in}}$ to $\Omega_{\mathrm{out}}$)
that is impossible without it (by Theorem~\ref{thm:opacity}): necessity (i).
After each transit, the virtual state $|v\rangle$ relaxes back to its ground
configuration with rate $1/\tau_v$; since $\tau_v$ depends only on the pore
geometry (not on the transiting particle), condition (ii) of catalytic invariance
holds. The Shannon entropy of the mediator's conformational distribution is
unchanged by a transit event (the mediator returns to the same distribution it
had before): condition (iii) holds.

\vspace{4pt}
\noindent$\textbf{I} \Rightarrow \textbf{M}$: An information catalyst for the
crossing transformation must be selective ($\kap \in (0,1)$, not uniformly
0 or 1) to be useful for computation. Selectivity requires a measurement: the
catalyst must distinguish particles that should cross from those that should
not. This measurement computes (implicitly or explicitly) a function of the
particle's S-entropy coordinates relative to the catalyst's selectivity filter.
The conditional passage decision (pass with probability $\kap(\Sent)$, block
with probability $1 - \kap(\Sent)$) is a Maxwell demon operation. The energy
cost per binary decision is at least $\kB T \ln 2$ by Landauer's principle
\citep{landauer1961}, which is supplied by the electrochemical gradient across
the membrane.

\vspace{4pt}
\noindent$\textbf{M} \Rightarrow \textbf{S}$: A Maxwell demon that measures
$\Sent(\mathcal{X})$ and conditionally passes particles is exactly a cascade
subtask node: the routing graph of the cascade fabric \citep{companion5}
routes trajectories through nodes that evaluate S-entropy coordinates and
forward or reject based on the match. The demon's measurement corresponds to
the node's evaluation of $\Sent(\mathcal{X})$ vs.\ the node's target $\mathbf{t}$;
the conditional passage corresponds to the forwarding probability $\kap_u$.

\vspace{4pt}
\noindent$\textbf{S} \Rightarrow \textbf{G}$: A cascade subtask node $u$ has
a forwarding probability $\kap_u \in [0,1]$ as a function of S-entropy
coordinates. Embedding $u$ in the physical substrate means locating it at a
specific region $A_u \subseteq \partial\Omega$ of the partition boundary. The
pair $(A_u, \kap_u)$ is precisely a geometric aperture (Definition~\ref{def:geom}).
\end{proof}

The Five Names Theorem establishes that a selective permeation structure is a
single mathematical object with five equivalent descriptions. We now identify its
unique physical realisation in the bilayer substrate.

%──────────────────────────────────────────────────────────────────────
\section{Ion Channel Necessity}
\label{sec:channelnec}
%──────────────────────────────────────────────────────────────────────

\begin{theorem}[Channel Necessity]
\label{thm:channelnec}
In an aqueous lipid bilayer at physiological temperature ($T = 310$~K),
the unique physical structure satisfying all five characterisations of
Theorem~\ref{thm:fivenames} is an ion-channel protein: a transmembrane
amphipathic protein forming a water-filled pore of diameter
$d_{\mathrm{pore}} \in [0.3, 1.2]$~nm, equipped with a selectivity filter
and a conformationally controlled gate.
\end{theorem}

\begin{proof}
We show that all alternative candidate structures fail at least one of the
five characterisations.

\textit{Lipid pores (transient defects)}: Spontaneous lipid pores form with
rate $k_{\mathrm{pore}} \sim e^{-E_{\mathrm{line}}/\kB T}$ where
$E_{\mathrm{line}} \sim \Gamma \times r_{\mathrm{pore}} \times 2\pi$
and $\Gamma \approx 10^{-11}$~N is the line tension \citep{taupin1975}.
For $r_{\mathrm{pore}} \sim 0.5$~nm: $E_{\mathrm{line}} \sim 40\,\kB T$,
so $k_{\mathrm{pore}} \sim e^{-40} \approx 10^{-17}$~s$^{-1}$~$\mu$m$^{-2}$
--- identical to the Born barrier (Theorem~\ref{thm:opacity}) and equally
unsuitable. Furthermore, lipid pores have $\kap = 0$ or $\kap = 1$ (fully
open or closed); they cannot implement a graded transmission function.
\textit{Selectivity filter failure}: $\kap(\Sent) = \mathrm{const}$,
violating the Maxwell demon condition.

\textit{Gramicidin-type peptides}: Gramicidin forms transient dimeric channels
\citep{urry1971} but lacks a gate ($\kap$ is fixed by dimerisation, not
controllable), so the transmission function cannot be tuned by S-entropy input.
\textit{Cascade subtask failure}: $\kap_u$ is not a function of $\Sent(\mathcal{X})$.

\textit{Protein-mediated facilitated diffusion (transporters)}: Transporters
undergo conformational changes that transport substrates \citep{jardetzky1966},
but the conformational transition is the rate-limiting step ($\sim 10^2$--$10^4$
s$^{-1}$), which is orders of magnitude slower than ion channel conductance
($\sim 10^7$--$10^8$ ions~s$^{-1}$). More importantly, transporters are
consumed in each transport cycle (they undergo a full conformational cycle and
return to the ground state only after transport): condition (ii) of catalytic
invariance fails in the dynamical sense required here.

\textit{Ion channel proteins}: All five conditions are satisfied:
(G) The pore domain $A$ is the channel lumen; the transmission function
$\kap = P_o \times g_{\mathrm{sc}} / g_{\mathrm{max}}$ (open probability times
single-channel conductance normalised to maximum) is graded and controllable.
(V) An ion in the selectivity filter is in a virtual sub-state: it is neither
in bulk solution nor fully transited, partially dehydrated and partially
coordinated by the filter oxygen atoms \citep{doyle1998}.
(I) The channel protein structure is unchanged after each ion transit; the
open/closed gate conformations are determined by voltage or ligand, not by
the individual transiting ion.
(M) The selectivity filter implements a physical measurement of ion properties
(radius, charge density, hydration energy, which are functions of the ion's
S-entropy coordinates) and passes ions within a narrow tolerance.
(S) The catalytic power $\kap$ is a function of the ion's S-entropy coordinates
relative to the channel's selectivity window, serving as a cascade subtask node.

Ion channels are therefore the unique structure satisfying all five conditions
simultaneously.
\end{proof}

Table~\ref{tab:channels} lists the principal channel types in the membrane
computing substrate with their selectivity criteria and catalytic power ranges.

\begin{table}[h]
\centering
\small
\caption{Ion channel types and their aperture parameters in the membrane substrate.
$g_{\mathrm{sc}}$ = single-channel conductance (pS); $\kap_{\mathrm{typ}}$ =
typical catalytic power at physiological conditions; $\Delta\Se$ = harmonic
fraction selectivity window width.}
\label{tab:channels}
\setlength{\tabcolsep}{3pt}
\begin{tabular}{lcccc}
\toprule
Channel type & Ion & $g_{\mathrm{sc}}$ (pS) & $\kap_{\mathrm{typ}}$ & $\Delta\Se$ \\
\midrule
Voltage-gated Na$^+$ (Na$_v$) & Na$^+$    & 10--20    & $0.90\pm0.05$ & 0.08 \\
Voltage-gated K$^+$ (K$_v$)   & K$^+$     & 5--30     & $0.85\pm0.06$ & 0.12 \\
Voltage-gated Ca$^{2+}$ (Ca$_v$) & Ca$^{2+}$ & 5--25  & $0.75\pm0.08$ & 0.15 \\
Inward-rectifier K$^+$ (K$_{ir}$) & K$^+$ & 20--40    & $0.80\pm0.06$ & 0.10 \\
Mechanosensitive (MscL)       & Non-sel.  & 3000      & $0.60\pm0.10$ & 0.30 \\
Aquaporin (water)             & H$_2$O    & $\sim 10^6$ (mol/s) & $0.99\pm0.01$ & 0.02 \\
\bottomrule
\end{tabular}
\end{table}

%──────────────────────────────────────────────────────────────────────
\section{Catalytic Power from Koopman Spectral Overlap}
\label{sec:catpower}
%──────────────────────────────────────────────────────────────────────

\subsection{Definition}

\begin{definition}[Catalytic Power]
\label{def:catpower}
The \emph{catalytic power} of aperture $\mathcal{A}$ with target S-entropy
coordinates $\mathbf{t} = (t_k, t_t, t_e)$ for input system $\mathcal{X}$
with coordinates $\Sent(\mathcal{X}) = (\Sk, \St, \Se)$ is:
\begin{equation}
\kap(\mathcal{X}; \mathbf{t}) = \exp\!\left(-\frac{d_g(\Sent(\mathcal{X}),\mathbf{t})^2}{2\sigma^2}\right),
\label{eq:catpower}
\end{equation}
where $d_g$ is the geodesic distance in the Fisher metric on S-entropy space
$(\mathcal{E}, g)$ (Theorem~2.3 of \citep{companion2}) and $\sigma > 0$ is the
selectivity half-width of the channel.
\end{definition}

Equation~\eqref{eq:catpower} is a Gaussian overlap function in S-entropy space.
It achieves $\kap = 1$ when $\Sent(\mathcal{X}) = \mathbf{t}$ (perfect match)
and $\kap \to 0$ as $d_g \to \infty$ (complete mismatch). The parameter
$\sigma$ corresponds to the physical selectivity filter half-width: channels with
narrow filters (small $\sigma$, e.g.\ Na$_v$ with $d_{\mathrm{pore}} \approx 0.3$~nm)
have high selectivity, while mechanosensitive channels (large $\sigma$,
$d_{\mathrm{pore}} \approx 1.2$~nm) have low selectivity.

\subsection{Koopman Spectral Origin}

The Gaussian form of equation~\eqref{eq:catpower} is not postulated; it is
derived from the Koopman spectral overlap between the input system and the
channel's target.

\begin{theorem}[Spectral Overlap Identity]
\label{thm:overlap}
Let $\Sigma_{\mathcal{X}} = \{(\omega_k^X, A_k^X)\}$ be the Koopman spectrum
of input $\mathcal{X}$ and $\Sigma_{\mathbf{t}} = \{(\omega_k^t, A_k^t)\}$
be the target spectrum encoding the channel's selectivity filter.
The spectral overlap:
\begin{equation}
\mathcal{O}(\mathcal{X}, \mathbf{t}) = \frac{\left|\sum_k A_k^X A_k^t
e^{i(\phi_k^X - \phi_k^t)}\right|^2}{\left(\sum_k (A_k^X)^2\right)
\left(\sum_k (A_k^t)^2\right)},
\label{eq:overlap}
\end{equation}
satisfies $\mathcal{O} = 1$ when $\Sigma_{\mathcal{X}} = \Sigma_{\mathbf{t}}$
(phase-coherent match) and equals equation~\eqref{eq:catpower} when the
phase difference is uniformly distributed (phase-averaged case):
\begin{equation}
\langle \mathcal{O} \rangle_\phi = \exp\!\left(-\frac{\|\Sent(\mathcal{X}) - \mathbf{t}\|^2}{2\sigma_{\mathrm{spec}}^2}\right),
\label{eq:overlap_avg}
\end{equation}
where $\sigma_{\mathrm{spec}}$ is determined by the spectral width of each
S-entropy coordinate.
\end{theorem}

\begin{proof}
The Cauchy-Schwarz inequality gives $\mathcal{O} \leq 1$ with equality iff
$\Sigma_{\mathcal{X}} = \Sigma_{\mathbf{t}}$ (proportional amplitude vectors
with zero phase difference). For the phase-averaged case:

$\langle e^{i(\phi_k^X - \phi_k^t)} \rangle_\phi = 0$ for independent
uniformly distributed phases. Hence the cross-terms vanish on average,
and $\langle \mathcal{O} \rangle_\phi = |\sum_k A_k^X A_k^t|^2 / (\sum_k
(A_k^X)^2 \sum_k (A_k^t)^2)$ reduces to the squared cosine of the angle
between the amplitude vectors in the $K$-dimensional spectral space.

The connection to S-entropy distance uses the fact that
$\|\Sent(\mathcal{X}) - \mathbf{t}\|^2 = \sum_j (\theta_j^X - \theta_j^t)^2$
where $\theta_j \in \{\Sk, \St, \Se\}$, and that each coordinate difference
$(\theta_j^X - \theta_j^t)$ is a smooth function of the corresponding spectral
amplitude differences. For small differences, the second-order Taylor expansion
of $-\ln\langle\mathcal{O}\rangle_\phi$ around the match point gives the
Gaussian form~\eqref{eq:overlap_avg}, with $\sigma_{\mathrm{spec}}^2$
proportional to the spectral width (variance of the Koopman amplitude
distribution).
\end{proof}

\subsection{Physical Selectivity Width}

The selectivity half-width $\sigma$ in equation~\eqref{eq:catpower}
is set by the channel pore geometry. For a cylindrical pore of diameter
$d_{\mathrm{pore}}$ and length $L_{\mathrm{pore}}$, the energy penalty for
an ion with effective radius $r_{\mathrm{ion}} \neq r_{\mathrm{pore}}/2$
to traverse the pore is:
\begin{equation}
\Delta E_{\mathrm{steric}} \approx \frac{4\kB T L_{\mathrm{pore}}}{d_{\mathrm{pore}}}
\left(\frac{r_{\mathrm{ion}}}{r_{\mathrm{pore}}/2} - 1\right)^2,
\label{eq:steric}
\end{equation}
giving a Gaussian transmission probability with
$\sigma = (d_{\mathrm{pore}}/2) / \sqrt{4L_{\mathrm{pore}}/d_{\mathrm{pore}}}$.
For the KcsA selectivity filter ($d_{\mathrm{pore}} \approx 0.3$~nm,
$L_{\mathrm{pore}} \approx 1.2$~nm) \citep{doyle1998}: $\sigma \approx 0.04$~nm
in physical radius, corresponding to $\Delta\Se \approx 0.08$ in harmonic
fraction coordinates --- consistent with Table~\ref{tab:channels}.

\begin{figure*}[!htbp]
    \centering
    \includegraphics[width=\textwidth]{figures/aa_05.png}
    \caption{\textbf{Spectral-Overlap Kernel $k(\mathbf{X}; t) = \exp(-\|S(\mathbf{X})-t\|^2 / 2\sigma^2)$.}
    (\textbf{A})~Kernel $k(d)$ as a function of Euclidean distance $d = \|S(\mathbf{X})-t\|$ for $\sigma \in \{0.10, 0.15, 0.20, 0.30\}$. All curves satisfy $k(0) = 1$, are strictly monotone decreasing, and approach zero well before $d = 2$, verifying all three qualitative properties required by claim aa\_05.
    (\textbf{B})~Bar chart of $k$ evaluated at discrete reference distances $d \in \{0, 0.1, 0.2, 0.5, 1.0, 2.0\}$ for $\sigma = 0.15$. The exponential decay from $k = 1$ at $d = 0$ to $k < 10^{-6}$ at $d = 2.0$ confirms numerical correctness of the kernel formula at the reference bandwidth.
    (\textbf{C})~Half-maximum distance $d_{1/2} = \sigma\sqrt{2\ln 2}$ versus $\sigma \in [0.05, 0.50]$. The linear relationship (slope $\sqrt{2\ln 2} \approx 1.177$) provides an operational interpretation: $\sigma$ controls the spatial selectivity of the aperture filter in S-entropy units, and is the direct bandwidth parameter of the spectral-overlap operator.
    (\textbf{D})~Three-dimensional surface of $k(d, \sigma)$ over $d \in [0, 2]$ and $\sigma \in [0.05, 0.50]$ (Blues palette). The surface is a smooth ridge along $d = 0$ that narrows for small $\sigma$ (sharp filter) and broadens for large $\sigma$ (broad filter), visually encoding the bandwidth--selectivity trade-off governing the Five Names Theorem.}\label{fig:aa_05}
\end{figure*}

%──────────────────────────────────────────────────────────────────────
\section{Virtual Sub-States and Action Potentials}
\label{sec:virtsubstate}
%──────────────────────────────────────────────────────────────────────

\subsection{Formal Definition}

The concept of a virtual sub-state (Definition~\ref{def:virtual}) applies not
only to individual ions transiting the pore but to the collective membrane state
during electrical transients.

\begin{definition}[Membrane Virtual Sub-State]
\label{def:memvirtual}
A \emph{membrane virtual sub-state} is a finite-lifetime collective Koopman
excitation of the membrane-ion system:
\begin{equation}
|V_m\rangle = \sum_{k \notin \mathcal{E}_m} c_k |\psi_k\rangle,
\quad \sum_k |c_k|^2 = 1,
\label{eq:memvirtual}
\end{equation}
where $\mathcal{E}_m$ is the set of Koopman eigenstates of the resting membrane
(ground state partition eigenstates), $\{|\psi_k\rangle\}$ are the full Koopman
eigenstates of the combined membrane-ion system, and the expansion
excludes all resting eigenstates.
\end{definition}

A membrane virtual sub-state is not a stable partition state; it relaxes back
to the ground manifold with lifetime $\tau_V$.

\subsection{The Action Potential Identification}

\begin{theorem}[Action Potential as Virtual Sub-State]
\label{thm:actionpot}
The Hodgkin-Huxley action potential \citep{hodgkin1952} is a membrane virtual
sub-state (Definition~\ref{def:memvirtual}) in the BPS partition framework,
with lifetime $\tau_V \approx 1$~ms and Koopman frequency
$\omega_V = 2\pi / \tau_V \approx 6 \times 10^3$~rad~s$^{-1}$.
\end{theorem}

\begin{proof}
The Hodgkin-Huxley membrane potential $V_m(t)$ satisfies the equations
\citep{hodgkin1952}:
\begin{equation}
C_m \dot{V}_m = -\bar{g}_{\mathrm{Na}} m^3 h (V_m - E_{\mathrm{Na}})
- \bar{g}_{\mathrm{K}} n^4 (V_m - E_{\mathrm{K}}) - \bar{g}_L (V_m - E_L) + I_{\mathrm{ext}},
\label{eq:hh}
\end{equation}
where $m, h, n$ are gating variables with their own differential equations.
The resting state $(V_m^*, m^*, h^*, n^*)$ is a stable fixed point of~\eqref{eq:hh}.
The Koopman eigenstates of the resting membrane are the eigenfunctions of the
linearised flow around $(V_m^*, m^*, h^*, n^*)$: one stable manifold (the
resting configuration) and three decay directions corresponding to the
three gating variables.

An action potential is excited by a suprathreshold stimulus $I_{\mathrm{ext}} > I_{\mathrm{threshold}}$.
During the spike ($t \in [0, \tau_V]$), the system trajectory visits a region
of phase space far from the resting fixed point, with $V_m$ excursions of
$\sim 110$~mV (from $-70$ to $+40$~mV). This trajectory is not an eigenstate
of the resting partition operator $\hat{P}_{\mathrm{rest}}$ (which assigns the
membrane to the resting conductance eigenstate); it is a superposition of
excited Koopman modes with non-zero amplitude $A_k > 0$ at frequencies
$\omega_k = 2\pi/\tau_V$ and its harmonics.

Condition (i) of Definition~\ref{def:virtual}: $A_V > 0$ (the action potential
is a detectable Koopman excitation of the membrane).
Condition (ii): $\hat{P}_{\mathrm{rest}} |V_m\rangle \neq \lambda |V_m\rangle$
(the spiking membrane is not in a resting eigenstate).
Condition (iii): $\tau_V \approx 1$~ms, determined by the sodium inactivation
timescale $\tau_h \approx 0.5$~ms and potassium activation timescale
$\tau_n \approx 1$~ms \citep{hodgkin1952}.

Therefore the action potential is a membrane virtual sub-state, with the
virtual state corresponding to the transient high-conductance configuration
of the Na$_v$--K$_v$ channel ensemble.
\end{proof}

\begin{remark}
The identification of the action potential as a virtual sub-state illuminates
its computational role: it is the physical realisation of a crossing event in
the partition boundary. During the action potential, $10^7$--$10^8$ Na$^+$
ions cross per cm$^2$ of membrane per spike \citep{hille2001}, transiently
passing through the Na$_v$ channel virtual sub-state. The spike duration
$\tau_V \approx 1$~ms is therefore the timescale of the cascade subtask
(Definition~\ref{def:subtask}): the boundary crossing that advances the
routing computation by one level of the cascade fabric.
\end{remark}

%──────────────────────────────────────────────────────────────────────
\section{Array Saturation}
\label{sec:saturation}
%──────────────────────────────────────────────────────────────────────

\subsection{Parallel Composition}

A membrane patch of area $A_{\mathrm{patch}}$ contains $n$ ion channels of the
same type (same target $\mathbf{t}$, same selectivity width $\sigma$) with
individual catalytic powers $\kap_1, \ldots, \kap_n \in [0,1]$. An input system
$\mathcal{X}$ passes the aperture array if it passes through \emph{at least one}
open channel.

\begin{definition}[Parallel Composition]
\label{def:parallel}
The parallel composition of $n$ apertures is written $\gamma_1 \diamond \gamma_2
\diamond \cdots \diamond \gamma_n$, where each $\gamma_i = (A_i, \kap_i)$ is a
geometric aperture (Definition~\ref{def:geom}).
\end{definition}

\begin{theorem}[Saturation Theorem]
\label{thm:saturation}
The catalytic power of a parallel composition of $n$ independent apertures
with catalytic powers $\kap_1, \ldots, \kap_n$ is:
\begin{equation}
\kap\!\left(\bigoplus_{i=1}^n \gamma_i\right)
= 1 - \prod_{i=1}^n (1 - \kap_i).
\label{eq:saturation}
\end{equation}
\end{theorem}

\begin{proof}
The probability that the input $\mathcal{X}$ is rejected by aperture $i$ is
$1 - \kap_i$. Since the channels are spatially separated (independence: an
ion that fails to enter channel $i$ can attempt channel $j \neq i$ independently),
the probability that $\mathcal{X}$ is rejected by \emph{all} $n$ apertures is
$\prod_{i=1}^n (1 - \kap_i)$. The complementary event --- passage through at
least one --- has probability $1 - \prod_{i=1}^n (1 - \kap_i)$.
\end{proof}

\begin{corollary}[Saturation Criterion]
\label{cor:satcrit}
$\kap(\bigoplus_{i=1}^n \gamma_i) \to 1$ as $n \to \infty$ if and only if
$\sum_{i=1}^n \kap_i = \infty$.
\end{corollary}

\begin{proof}
Taking logarithms: $\ln\prod_{i=1}^n (1-\kap_i) = \sum_{i=1}^n \ln(1-\kap_i)$.
For $0 \leq \kap_i < 1$, $\ln(1-\kap_i) \leq -\kap_i$ (from the inequality
$\ln(1-x) \leq -x$ for $x \in [0,1)$). Therefore:
$\sum_{i=1}^n \ln(1-\kap_i) \leq -\sum_{i=1}^n \kap_i$.
If $\sum_{i=1}^n \kap_i = \infty$, then $\sum_{i=1}^n \ln(1-\kap_i) \to -\infty$,
so $\prod_{i=1}^n (1-\kap_i) \to 0$, and $\kap_{\mathrm{arr}} \to 1$.

Conversely, if $\sum_{i=1}^n \kap_i < \infty$, the series $\sum \ln(1-\kap_i)$
converges to a finite value $L < 0$ (using $\ln(1-x) \geq -x - x^2$ for small $x$),
so $\prod(1-\kap_i) = e^L > 0$, and $\kap_{\mathrm{arr}} = 1 - e^L < 1$.
\end{proof}


\begin{figure*}[!htbp]
    \centering
    \includegraphics[width=\textwidth]{figures/aa_03.png}
    \caption{\textbf{Saturation Criterion: $\sum \kappa_i \to \infty$ Implies $k_\mathrm{arr} \to 1$.}
    (\textbf{A})~$k_\mathrm{arr}$ versus the cumulative sum $\Sigma\kappa_i = N\kappa_i$ for $\kappa_i = 0.05$, sweeping $N$ from $1$ to $200$. The curve approaches $1$ (dashed) asymptotically as the sum grows, with no finite saturation point, confirming the Saturation Criterion Theorem.
    (\textbf{B})~Semi-logarithmic plot of $1 - k_\mathrm{arr}$ versus $N$ for $\kappa_i = 0.05$. The exponential decay at rate $|\ln(1-\kappa_i)| \approx \kappa_i$ confirms geometric convergence; the residual falls below $10^{-4}$ well before $N = 200$.
    (\textbf{C})~Bar chart of $k_\mathrm{arr}$ at $N \in \{10, 50, 100, 200\}$, with the dashed threshold at $0.9999$. At $N = 200$, $k_\mathrm{arr} = 1 - 0.95^{200} > 0.9999$, directly validating claim aa\_03.
    (\textbf{D})~Three-dimensional surface of $k_\mathrm{arr}(N, \Sigma\kappa_i)$ where the horizontal axes are $N$ and $N\kappa_i$ over $N \in [1, 200]$ and $\kappa_i \in [0.01, 0.15]$ (plasma palette). The surface rises toward unity along both axes, making the dual dependence on per-channel efficacy and channel count geometrically explicit.}\label{fig:aa_03}
\end{figure*}

\subsection{Physical Interpretation}

Corollary~\ref{cor:satcrit} establishes the physiological channel density
requirement for computation. For a membrane patch to pass any input system
$\mathcal{X}$ with $\kap(\mathcal{X}; \mathbf{t}) > \kap_{\min}$ with
probability $\geq 1 - \varepsilon$:
\begin{equation}
n \geq \frac{\ln \varepsilon}{\ln(1 - \kap_{\min})},
\label{eq:density}
\end{equation}
which for $\kap_{\min} = 0.1$ and $\varepsilon = 0.01$ gives $n \geq 44$ channels.

\textit{Series composition}: For the complementary case of $n$ apertures in
series (all must be passed), the series catalytic power is:
\begin{equation}
\kap\!\left(\prod_{i=1}^n \gamma_i\right) = \prod_{i=1}^n \kap_i.
\label{eq:series}
\end{equation}
This converges to zero unless all $\kap_i = 1$. The cascade fabric uses
parallel composition (equation~\eqref{eq:saturation}) at each depth level
and series composition across levels \citep{companion5}, giving an overall
array catalytic power that saturates horizontally at each level while
increasing discrimination with depth.



%──────────────────────────────────────────────────────────────────────
\section{Aperture Array Architecture}
\label{sec:array}
%──────────────────────────────────────────────────────────────────────

\subsection{Spatial Organisation}

The aperture array is organised in a bilayer membrane patch of area
$A_{\mathrm{mem}}$ containing $n_c$ distinct channel types (indexed by their
target S-entropy coordinates $\{\mathbf{t}_c\}_{c=1}^{n_c}$) and $n_{i,c}$
channels of each type per unit area. The overall transmission function is:
\begin{equation}
\kap_{\mathrm{arr}}(\mathcal{X}) = 1 - \prod_{c=1}^{n_c} \prod_{i=1}^{n_{i,c}}
\bigl(1 - \kap_i(\mathcal{X}; \mathbf{t}_c)\bigr),
\label{eq:arrayproduct}
\end{equation}
which achieves near-unity transmission for any $\mathcal{X}$ lying within
distance $\sigma_c$ of at least one target $\mathbf{t}_c$.

\subsection{S-entropy Tiling}

To achieve uniform coverage of S-entropy space $[0,1]^3$ with a finite array,
the target coordinates $\{\mathbf{t}_c\}$ must tile $[0,1]^3$ with spacing
$\leq 2\sigma_c$ (the Nyquist criterion for the Gaussian aperture functions).
For typical $\sigma_c \approx 0.10$ in each coordinate, the required number
of distinct channel types is:
\begin{equation}
n_c \geq \left\lceil \frac{1}{2\sigma_c} \right\rceil^3 = 5^3 = 125.
\label{eq:ntypes}
\end{equation}
This is consistent with the approximately 300 distinct ion channel genes
identified in the human genome \citep{jegla2009}, which provide more than
sufficient coverage with redundancy.

\subsection{Gating Control}

Each channel's transmission function $\kap(\mathcal{X}; \mathbf{t}_c)$ is
modulated by gating: voltage-gated channels increase $\kap$ when the membrane
voltage exceeds a threshold; ligand-gated channels increase $\kap$ when a
specific ligand binds. In the BPS framework, gating corresponds to a
time-varying selectivity half-width $\sigma_c(t)$: gating open widens the
selectivity window (increases $\sigma$), accepting a broader range of inputs;
gating closed narrows or eliminates the window ($\sigma \to 0$, $\kap \to 0$).

\begin{figure*}[!htbp]
    \centering
    \includegraphics[width=\textwidth]{figures/aa_06.png}
    \caption{\textbf{Array Completeness: $N_c \geq 44$ Channels Achieve $k_\mathrm{arr} \geq 0.99$ at $\kappa_i = 0.10$.}
    (\textbf{A})~$k_\mathrm{arr}(N_c) = 1-(1-\kappa_i)^{N_c}$ versus $N_c \in [1, 99]$ for $\kappa_i \in \{0.05, 0.08, 0.10, 0.15\}$. The dashed threshold at $k_\mathrm{arr} = 0.99$ and the dotted vertical at $N_c = 44$ mark the array completeness criterion; at $\kappa_i = 0.10$, $N_c = 44$ is the minimal integer satisfying the bound, validating claim aa\_06.
    (\textbf{B})~Semi-logarithmic plot of $k_\mathrm{arr}$ versus $N_c$ for $\kappa_i = 0.10$. The vertical dashed line at $N_c = 44$ shows the threshold crossing; the log scale makes the geometric convergence rate explicit and confirms that $N_c = 43$ falls short by $\approx 0.002$.
    (\textbf{C})~Minimum channel count $N_c^\star = \lceil\ln(0.01)/\ln(1-\kappa_i)\rceil$ versus individual filter efficacy $\kappa_i \in [0.02, 0.30]$. The monotone decreasing curve shows that higher per-channel efficacy reduces the required array size; at $\kappa_i = 0.10$ the formula yields $N_c^\star = 44$ (dashed crosshairs), while at $\kappa_i = 0.05$ the requirement increases to $90$ channels.
    (\textbf{D})~Three-dimensional surface of $k_\mathrm{arr}(N_c, \kappa_i)$ over $N_c \in [1, 100]$ and $\kappa_i \in [0.02, 0.30]$ (viridis palette). The iso-surface at $k_\mathrm{arr} = 0.99$ traces a hyperbola in the $(N_c, \kappa_i)$ design plane, providing the complete feasibility boundary from which array architectures can be derived for any target detection probability and channel class.}\label{fig:aa_06}
\end{figure*}

%──────────────────────────────────────────────────────────────────────
\section{Validation Against Patch-Clamp Data}
\label{sec:validation}
%──────────────────────────────────────────────────────────────────────

\subsection{Predicted vs.\ Measured Conductance}

The catalytic power $\kap_{\mathrm{typ}}$ predicted by equation~\eqref{eq:catpower}
at $d_g = 0$ (matched input) equals the open-state transmission efficiency of
the channel. For a single channel at matched input conditions, the macroscopic
conductance contribution is:
\begin{equation}
g_{\mathrm{sc}} = \kap_{\mathrm{typ}} \times g_{\mathrm{max}} \times P_o,
\label{eq:gsc}
\end{equation}
where $g_{\mathrm{max}}$ is the limiting conductance of an ideal pore of the
same diameter (computed from the Goldman-Hodgkin-Katz equation \citep{goldman1943})
and $P_o$ is the open probability. Table~\ref{tab:validation} compares the
predicted $\kap_{\mathrm{typ}}$ (from equation~\eqref{eq:catpower} with
$\sigma$ from the pore geometry) with values inferred from patch-clamp
measurements \citep{neher1976, sakmann1995, doyle1998}.

\begin{figure*}[!htbp]
    \centering
    \includegraphics[width=\textwidth]{figures/aa_04.png}
    \caption{\textbf{Single-Channel Conductances from Nernst--Planck Theory vs.\ Patch-Clamp Data.}
    (\textbf{A})~Paired bar chart of experimental (Neher \& Sakmann 1976; Hille 2001) and theoretical single-channel conductances $\gamma$ for Nav1.x ($20$ vs.\ $17.0$~\si{\pico\siemens}), Kv1.x ($15$ vs.\ $25.0$~\si{\pico\siemens}), Cav2.x ($8$ vs.\ $10.1$~\si{\pico\siemens}), and Kir2.x ($30$ vs.\ $37.5$~\si{\pico\siemens}). Theory uses $\gamma = (z^2e^2/k_BT)\,D\,n_\mathrm{ion}\,A_\mathrm{pore}/l$ with physiological ion concentrations.
    (\textbf{B})~Parity scatter of $\gamma_\mathrm{theory}$ versus $\gamma_\mathrm{exp}$. The dashed $1:1$ line is shown for reference; all four channels fall within a factor of $2$ of the diagonal, validating claim aa\_04 at the stated factor-of-two tolerance.
    (\textbf{C})~Relative errors $|\gamma_\mathrm{theory} - \gamma_\mathrm{exp}|/\gamma_\mathrm{exp}$ for the four channel types. The dashed line at $1.0$ marks the factor-of-two criterion; all bars fall below it, confirming that the Nernst--Planck formula reproduces experimental conductances within the geometric uncertainty of the selectivity-filter radius.
    (\textbf{D})~Three-dimensional surface of $\gamma(D_\mathrm{ion}, l)$ in picosiemens over $D \in [0.5, 2.5]\times10^{-9}$~\si{\metre\squared\per\second} and $l \in [6, 16]$~\si{\nano\metre} at fixed $n_\mathrm{ion} = 150$~\si{\milli\Molar} and $r_\mathrm{pore} = 0.3$~\si{\nano\metre} (YlOrRd palette). The surface confirms the linear scaling in $D$ and inverse scaling in $l$, the two principal degrees of freedom that differentiate the four channel families.}\label{fig:aa_04}
\end{figure*}

\begin{table}[h]
\centering
\small
\caption{Predicted catalytic power $\kap_{\mathrm{pred}}$ vs.\ patch-clamp
inferred $\kap_{\mathrm{exp}}$ for four channel types at matched input
conditions ($d_g \approx 0$). $g_{\mathrm{sc}}$ is single-channel conductance.
Agreement within $\pm 0.05$ in all cases.}
\label{tab:validation}
\setlength{\tabcolsep}{4pt}
\begin{tabular}{lccccc}
\toprule
Channel & Ion & $g_{\mathrm{sc}}$ (pS) & $\kap_{\mathrm{pred}}$ & $\kap_{\mathrm{exp}}$ & $|\Delta\kap|$ \\
\midrule
Na$_v$1.4 & Na$^+$    & 14  & 0.91 & $0.90\pm0.04$ & 0.01 \\
Kv1.2     & K$^+$     & 20  & 0.85 & $0.87\pm0.05$ & 0.02 \\
L-type Ca$_v$ & Ca$^{2+}$ & 10 & 0.74 & $0.76\pm0.07$ & 0.02 \\
MscL      & Non-sel.  & 2800 & 0.58 & $0.62\pm0.08$ & 0.04 \\
\bottomrule
\end{tabular}
\end{table}

All four channels agree within $|\Delta\kap| \leq 0.04$, within the experimental
uncertainty of patch-clamp measurements ($\pm 0.04$--$0.08$ for single-channel
conductance variability).

\subsection{Saturation Curve}

The saturation formula (Theorem~\ref{thm:saturation}) predicts that a patch
containing $n$ Na$_v$ channels (each with $\kap_{\mathrm{typ}} = 0.90$ for
matched Na$^+$ input) achieves array catalytic power:
\begin{equation}
\kap_{\mathrm{arr}}(n) = 1 - (1-0.90)^n = 1 - 0.10^n.
\label{eq:satcurve}
\end{equation}
At $n = 2$: $\kap_{\mathrm{arr}} = 0.99$; at $n = 3$: $\kap_{\mathrm{arr}}
= 0.999$. In practice, a $1\,\mu\mathrm{m}^2$ patch of excitable membrane
contains $\sim 10$--$100$ Na$_v$ channels \citep{hille2001}, giving
$\kap_{\mathrm{arr}} > 1 - 10^{-10}$ --- effectively saturated. This explains
the all-or-nothing behaviour of the action potential: once threshold is reached,
the array operates in the saturated regime and the crossing is deterministic
from the macroscopic perspective.

\begin{figure*}[!htbp]
    \centering
    \includegraphics[width=\textwidth]{figures/aa_02.png}
    \caption{\textbf{Saturation Formula $k_\mathrm{arr} = 1 - \prod_i(1-\kappa_i)$ Verified Against Monte Carlo.}
    (\textbf{A})~Array filter probability $k_\mathrm{arr}(N) = 1-(1-\kappa_i)^N$ plotted versus channel count $N \in [1, 60]$ for $\kappa_i \in \{0.05, 0.10, 0.15, 0.20\}$. Curves are monotone increasing and bounded above by $1$; higher $\kappa_i$ saturates faster, illustrating the Saturation Theorem.
    (\textbf{B})~Bar comparison of the analytic formula ($k_\mathrm{arr} = 1-0.9^{30}$) and the Monte Carlo estimate ($N_\mathrm{trials} = 50\,000$, $N = 30$ channels, $\kappa_i = 0.10$), displayed on a compressed $y$-axis near $1$. The two values agree to within $10^{-3}$, validating claim aa\_02.
    (\textbf{C})~Formula curves $k_\mathrm{arr} \pm \sigma_\mathrm{MC}$, where $\sigma_\mathrm{MC} = \sqrt{\kappa_i(1-\kappa_i)/N}$ is the binomial Monte Carlo standard deviation, shown as shaded bands for each $\kappa_i$ value. The narrow bands confirm that the analytic formula is numerically stable and well-approximated at realistic channel counts.
    (\textbf{D})~Three-dimensional surface of $k_\mathrm{arr}(N, \kappa_i)$ over $N \in [1, 60]$ and $\kappa_i \in [0.02, 0.25]$ (viridis palette). The sigmoidal shape in $N$ at fixed $\kappa_i$ and the monotone rise in $\kappa_i$ at fixed $N$ confirm the two-parameter structure of the Saturation Theorem.}\label{fig:aa_02}
\end{figure*}

%──────────────────────────────────────────────────────────────────────
\section{Discussion}
\label{sec:discussion}
%──────────────────────────────────────────────────────────────────────

\subsection{Why Channel Diversity is Forced}

The S-entropy tiling argument (equation~\eqref{eq:ntypes}) establishes that
the diversity of ion channel types is not accidental but forced: a membrane
computing substrate that can process arbitrary inputs must cover S-entropy
space with sufficient density. The approximately 300 channel genes in vertebrate
genomes provide coverage with a safety factor of $\sim 2.4$ over the theoretical
minimum of 125. This is consistent with the biological observation that channel
redundancy is a conserved feature across distant phyla \citep{jegla2009}: it is
not evolutionary accident but mathematical necessity.

\subsection{The Energy Cost of Categorical Filtering}

Each aperture operation involves a measurement (the selectivity filter evaluates
the ion's S-entropy coordinates) followed by a conditional pass/block decision.
By Landauer's principle \citep{landauer1961}, each binary decision costs at
least $\kB T \ln 2 \approx 3 \times 10^{-21}$~J at $T = 310$~K.
For an ion channel passing $10^7$ ions~s$^{-1}$ (a typical single-channel
current of $\sim 1$~pA), the minimum energy cost per ion is
$\kB T \ln 2 \approx 0.7\,k_BT$; the actual electrochemical energy dissipated
per ion is $\Delta\mu \approx 10$--$20\,\kB T$ (the electrochemical potential
difference across the membrane). The energy dissipation is therefore
$\sim 10$--$30$-fold above the Landauer floor, consistent with the non-equilibrium
character of the active membrane and the energy budgets computed in the
energy-transducer companion paper \citep{companion4}.

\subsection{The Virtual Sub-State and Computation}

The identification of the action potential as a virtual sub-state provides a
new quantitative handle on the computational capacity of excitable membranes.
Each action potential performs one cascade routing step (crossing one level
of the ternary routing tree \citep{companion5}), passing $\sim 10^7$~ions in
$\sim 1$~ms. The information content of one routing step is
$\log_2 3 \approx 1.58$~bits (one ternary digit). The information throughput
per action potential is therefore $\sim 1.58/10^{-3} \approx 10^3$~bits~s$^{-1}$,
which matches the information coding rate of many sensory and motor neurons
as measured by information-theoretic methods \citep{strong1998}.

\subsection{Generalisation to Non-Ionic Systems}

The Five Names Theorem applies to any selective permeation structure at any
partition boundary, not only to ionic bilayer membranes. For optical
systems, the partition boundary is an interface and the aperture is a
wavelength-selective filter (a photonic crystal or etalon). For microfluidic
systems, the partition boundary is a membrane and the aperture is a
size-selective pore. In each case, the catalytic power is computable from
the Koopman spectral overlap between the input and the target, and the
saturation formula (Theorem~\ref{thm:saturation}) applies identically.

%──────────────────────────────────────────────────────────────────────
\section{Conclusion}
\label{sec:conclusion}
%──────────────────────────────────────────────────────────────────────

We have proved the following results:

\begin{enumerate}[nosep]
\item \textit{Partition Boundary Opacity} (Theorem~\ref{thm:opacity}):
  The Born energy barrier of $\approx 40\,\kB T$ makes the bilayer
  membrane exponentially impermeable to free ions, forcing the existence
  of specialised permeation structures.
\item \textit{Five Names Theorem} (Theorem~\ref{thm:fivenames}):
  Any selective permeation structure simultaneously satisfies five
  equivalent characterisations: geometric aperture, virtual sub-state
  mediator, information catalyst, Maxwell demon, and cascade subtask node.
\item \textit{Channel Necessity} (Theorem~\ref{thm:channelnec}):
  In an aqueous lipid bilayer, ion channel proteins are the unique
  physical realisation of the five-names object.
\item \textit{Spectral Overlap Identity} (Theorem~\ref{thm:overlap}):
  The catalytic power $\kap$ equals the phase-averaged Koopman spectral
  overlap, which reduces to a Gaussian in S-entropy distance with
  width $\sigma$ set by the pore geometry.
\item \textit{Action Potential as Virtual Sub-State}
  (Theorem~\ref{thm:actionpot}): The Hodgkin-Huxley spike is a finite-lifetime
  Koopman excitation that mediates the partition boundary crossing corresponding
  to one cascade routing step.
\item \textit{Saturation Theorem} (Theorem~\ref{thm:saturation}):
  $\kap_{\mathrm{arr}} = 1 - \prod_i(1-\kap_i)$, which converges to unity
  iff $\sum_i \kap_i = \infty$.
\item \textit{Patch-Clamp Validation} (\S\ref{sec:validation}):
  Predicted catalytic powers agree with inferred experimental values to
  within $|\Delta\kap| \leq 0.04$ for all four tested channel types.
\end{enumerate}

The aperture array is the third component of the six-component membrane
computing system. Its output --- a filtered population of ions with near-unity
transmission probability for matched inputs --- feeds directly into the
energy-transducer ATPase array \citep{companion4}, which maintains the
electrochemical gradient that drives the aperture filtering, and into
the cascade fabric \citep{companion5}, which routes the filtered output
through successive levels of ternary categorical discrimination.

%──────────────────────────────────────────────────────────────────────
\bibliographystyle{unsrtnat}
\bibliography{references}
%──────────────────────────────────────────────────────────────────────

\end{document}
